\section{Phase 1: Hybrid Deployment}


\textbf{Status: Active (since Lux genesis).}

In Phase 1, both classical and post-quantum signatures coexist. The security guarantee is: the system is secure if \emph{either} the classical or the post-quantum scheme is unbroken.

\subsection{Dual Signature Model}

Every validator registers four keys on the P-Chain:

\begin{center}
\begin{tabular}{llrl}
\toprule
\textbf{Key Type} & \textbf{Algorithm} & \textbf{Size} & \textbf{Purpose} \\
\midrule
EC key & secp256k1 & 33 B (compressed) & EVM address derivation \\
BLS key & BLS12-381 & 48 B & Per-block consensus aggregation \\
ML-DSA key & ML-DSA-65 & 1,952 B & Per-epoch PQ identity \\
Pulsar key & Module-LWE & 2,048 B & Per-epoch PQ anonymous threshold \\
\bottomrule
\end{tabular}
\end{center}

\subsection{Per-Block Verification (Classical)}

Each block contains a BLS aggregate signature. Verification:

$$\text{BLS.Verify}(\text{aggPK}, \text{blockHash}, \sigma_{BLS}) \stackrel{?}{=} \text{true}$$

This provides sub-millisecond consensus for real-time block production.

\subsection{Per-Epoch Verification (PQ)}

Each epoch (1,000 blocks = 1 second) contains a STARK proof compressing $N$ ML-DSA signatures and a Pulsar threshold signature:

$$\text{STARK.Verify}(\pi, \text{epochDigest}) \stackrel{?}{=} \text{true}$$

This provides post-quantum finality without per-block overhead.

\subsection{Hybrid Security Property}

\begin{theorem}[Hybrid Security]
If either (a) the BLS discrete logarithm problem is hard, or (b) the Module-LWE problem (ML-DSA) and Module-LWE problem (Pulsar) are hard, then the consensus mechanism is secure.
\end{theorem}

\begin{proof}
The BLS aggregate provides classical security per-block. The STARK-compressed ML-DSA + Pulsar provides post-quantum security per-epoch. An adversary must break both to forge a finalized epoch. Since the two are based on independent mathematical problems (discrete log vs. lattice), compromising one does not help with the other.
\end{proof}

